\section{Introduction}
\label{sec:intro}

Graph neural networks (GNNs) are increasingly used as vision backbones: an image is represented as a set of superpixels or patches connected by a spatial adjacency, and message passing replaces the regular convolutional grid~\cite{dwivedi2023,han2022vig,achanta2012}. This representation is attractive for irregular objects and for event-driven deployment, yet the ANN-to-SNN conversion theory that makes spiking inference cheap was developed almost entirely for row-stochastic convolutional operators~\cite{diehl2015,rueckauer2017,sengupta2019,deng2021,bu2022qcfs}. For a graph, the propagation operator is neither a grid nor necessarily row-stochastic, and its row sums depend on degree and edge weight. Whether a conversion recipe calibrated for convolutions remains valid on such operators is largely untested.

A second, more practical gap concerns energy. Spiking inference replaces multiply--accumulates (MACs) with accumulations (ACs); under a 45\,nm model the per-operation ratio is only \MACACRatio~\cite{horowitz2014}. Consequently a $5\times$ arithmetic-energy advantage requires the total number of synaptic operations to fall below the ANN's MAC count, not merely to replace some of them. Most conversion studies report accuracy--latency curves on CNNs or ImageNet-scale models~\cite{wu2024ftbc,bojkovic2025tpp,hao2025parallel,huang2025dcgs}, and spike counts on graphs are rarely converted into an explicit ANN-relative energy statement with two independent accounting conventions.

This paper studies a combination of conversion components that was originally assembled for graph node classification, and asks whether that combination transfers to \emph{graph-structured visual classification}. Our starting point is a two-layer GCN with an exact conversion-error identity $\Bh - B = \RAEW$, a closed-form forward bias $\beta^\star=\theta/(2T)$, a membrane-potential readout, and a margin-recovery bound. We do \emph{not} claim to invent these individual components: threshold balancing, zero-expected-error bias, membrane decoding, deterministic counting and two-phase neurons all have their own lineage~\cite{deng2021,bu2022qcfs,wu2024ftbc,dualphase2024,bojkovic2025tpp,ramesh2025pascal}. Our contribution is the \emph{re-identification and adaptation of the combination} on a new operator class: we show that a linear global pooling matrix preserves the identity exactly, that the closed-form bias becomes redundant and harmful once the spike train is deterministic, and that margin regularisation plus a narrow, low-timestep network reaches a stable arithmetic-energy advantage while keeping the conversion fidelity measurable.

Specifically, we make four contributions.

\begin{itemize}
\item \textbf{Exact identity with a graph-level readout.} For a block-diagonal superpixel adjacency $A$ and a linear pooling matrix $R$, the conversion error is exactly $\Bh-B = R A (E_1 W_2)$ with $E_1 = c_1/T - \operatorname{ReLU}(B_1)$. We verify this identity to \IdentityMax\ with a unified float64 forward pass, and show that the residual observed under float32 is an accumulation artefact rather than a theoretical failure.
\item \textbf{A mechanism-level correction of the combination.} Deterministic two-phase spiking already realises the target firing count exactly ($\sum_t s[t]=n$). Adding $\beta^\star=\theta/(2T)$, which compensates floor-quantisation bias, double-corrects the rate. Removing it reduces the $T{=}8$ conversion loss from \BetaLossVoneM\ to \BetaLossVtwoM\ on MNIST superpixels and from \BetaLossVoneC\ to \BetaLossVtwoC\ on CIFAR10 superpixels, across five seeds.
\item \textbf{A stable arithmetic-energy result.} Across four mutually disjoint holdout groups (\TotalSeeds\ independent seeds), the converted network exceeds $5\times$ under both energy conventions on both datasets, without a single seed-level exception. Section~\ref{sec:experiments} reports the ratios and the margins over every adapted single-point baseline.
\item \textbf{An explicit accuracy--energy frontier.} Rather than optimising a single operating point, we characterise how conversion loss trades against energy and include a negative result: a larger training budget worsens fidelity at low timestep counts. The frontier, not a post-hoc operating point, is the object of study.
\end{itemize}

Energy numbers are 45\,nm \emph{theoretical arithmetic op-count estimates}, not measured joules on neuromorphic hardware. Because the energy constraint forces a narrow network, teacher accuracies (\MNISTAnn\% and \CIFARAnn\%) are modest; the evaluation therefore concerns conversion fidelity and energy rather than classifier quality. The method covers fixed, non-negative, linear adjacency with linear pooling; dynamic attention, nonlinear pooling and ImageNet-scale training fall outside its scope.
